\section{Motivation}

\newacronym{cern}{CERN}{European Organization for Nuclear Research}
\newacronym{sm}{SM}{Standard Model}
\newacronym{lhc}{LHC}{Large Hadron Collider}
\newacronym{lum}{$\mathcal{L}$}{luminosity}

In the \acrfull{lhc}, at \acrshort{cern}, the quest to understand the universe at its most fundamental level drives the core
of particle physics research. A significant part of this endeavor is centered around validating the \acrfull{sm}, a comprehensive
theory that explains the fundamental particles and forces in the universe. Particle colliders, key instruments in this research
process, accelerate and collide particles to reveal new insights through the analysis the byproducts of such collisions. The efficacy
of these colliders, relies heavily on a critical factor: the \acrfull{lum}. This parameter, quantifying the rate at which particles pass
through a specified area, directly influences the number of collisions a collider can produce. Furthermore, since the calculated cross
sections of particle interactions — a cornerstone in verifying Standard Model predictions — are highly dependent on precise luminosity
measurements, the accurate determination of luminosity becomes crucial in particle physics research.
